{\Large
	\noindent
	{\bf Zusammenfassung}} \\
\vspace*{0.3 cm}

\noindent
Die zunehmende Abhängigkeit moderner Elektrizitätssysteme von wetterabhängiger erneuerbarer Erzeugung hat die Anomalieerkennung in Strommärkten und im Netzbetrieb sowohl an Bedeutung gewinnen lassen als auch deren Komplexität erhöht. Traditionelle statistische Verfahren stoßen dabei häufig an Grenzen, wenn es darum geht, tatsächlich abnormes Verhalten von kontextbedingten Variationen zu unterscheiden, die durch meteorologische Bedingungen, saisonale Effekte und Marktvolatilität hervorgerufen werden. Die vorliegende Arbeit begegnet dieser Herausforderung durch die Entwicklung eines leichtgewichtigen neuro-symbolischen Rahmens für die Zeitreihen-basierte Anomalieerkennung im österreichischen Elektrizitätssektor. Der vorgeschlagene Ansatz kombiniert ein kompaktes Vorhersagemodell auf Basis diskreter Cellular Neural Networks und Extreme Learning Machines (dCeNN--ELM) mit Answer Set Programming (ASP) zur symbolischen Verfeinerung von Anomaliekandidaten.

Das vorgeschlagene System basiert auf einem aufbereiteten österreichischen Stunden-Datensatz für den Zeitraum 2017--2022 und integriert Strommarktvariablen, netzbetriebliche Signale, Wetterdaten, Informationen zur installierten Leistung sowie Kalenderkontext. Die neuronale Komponente modelliert das erwartete multivariate Systemverhalten und erzeugt residualbasierte Anomaliekandidaten, während die symbolische Komponente Persistenz-, Plausibilitäts- und Domänenkonsistenzregeln anwendet, um schwache oder unplausible Detektionen zurückzuweisen. Auf diese Weise geht der vorgeschlagene Rahmen über ein rein statistisches Anomalie-Scoring hinaus und ermöglicht eine besser interpretierbare sowie betrieblich aussagekräftigere Anomalieanalyse.

Die empirische Evaluierung umfasst die Vorhersagegüte, die Verfeinerung von Anomalien, eine finanzorientierte Nutzenbewertung sowie ereignisbasierte Fallstudien. Die Ergebnisse zeigen, dass der vorgeschlagene Rahmen die Baselines nicht in jeder Hinsicht hinsichtlich der reinen Vorhersagegenauigkeit übertrifft; insbesondere erweist sich Ridge Regression in den im Repository dokumentierten Ergebnissen als starker Referenzwert. Das dCeNN--ELM-Modell weist jedoch im Vergleich zur LSTM-Baseline deutlich geringere Trainings- und Inferenzkosten auf und ist stärker mit den Edge-orientierten sowie modularen Entwurfszielen dieser Arbeit vereinbar. Von besonderer Bedeutung ist darüber hinaus, dass die ASP-Verfeinerungsstufe die Interpretierbarkeit erhöht, indem sie eine explizite Veto-Logik für zurückgewiesene Anomalien bereitstellt, und durch domäneninformiertes Schließen zur Reduktion von Fehlalarmen beiträgt. Die ereignisbasierte Analyse verdeutlicht zudem, dass der vorgeschlagene Rahmen in der Lage ist, anhaltende Preisstörungen, markante Lastabweichungen, erneuerbarenbezogene Stressereignisse sowie mehrsignalige Anomalieereignisse zu identifizieren.

Insgesamt zeigt die vorliegende Arbeit, dass ein leichtgewichtiger neuro-symbolischer Ansatz eine tragfähige und praktisch nutzbare Alternative zur rein statistischen Anomalieerkennung in wetterabhängigen Elektrizitätssystemen darstellen kann. Der zentrale Beitrag der Arbeit liegt nicht in der Behauptung einer universellen prognostischen Überlegenheit, sondern in der Demonstration, wie Lernen, symbolisches Schließen und nachgelagerte Nutzenanalyse zu einer kohärenten Pipeline für die Anomalieanalyse integriert werden können.